%-------------------------------------------------------------------------------
%	SECTION TITLE
%-------------------------------------------------------------------------------
% Keep the heading and first project entry together.
\Needspace{4cm}
\cvsection{项目经历}


%-------------------------------------------------------------------------------
%	CONTENT
%-------------------------------------------------------------------------------
\begin{cventries}

  %---------------------------------------------------------
  \cventry
  {SoC/FPGA 验证与调试} % Role
  {Wujian100 SoC FPGA 原型验证与跨层调试} % Project
  {课程项目} % Course/lab/team/personal project
  {2026 年 6 月} % Date(s)
  {
    \begin{cvitems} % Description(s) of goals/contributions/results
      \item {\textbf{完整原型验证流程}：完成 Wujian100 SoC 从 C 语言 CDK 交叉编译、ELF-to-MEM 内存镜像转换到 vSyn 综合、vCom 编译、PnR 比特流生成以及 vDbg FPGA 上板测试的完整原型验证流程。}
      \item {\textbf{片上波形采集分析}：针对片上高速信号难以直接观测的问题，在 FPGA 原型中配置硬件探针并按时钟周期采集总线信号，使用 GTKWave 对目标输出地址的写入进行分析，将软件执行结果与实际硬件总线行为对应起来。}
      \item {\textbf{跨层定位工具链 Bug}：排查上板后程序未产生预期输出、而持续触发 CPU 异常的问题；从异常波形反向追踪程序执行状态，并结合 ELF 与内存映射分析，最终定位到工具链中 Python 脚本对内存地址空洞的错误处理，导致部分镜像整体向低地址偏移 4 Bytes，触发 CPU 非对齐加载异常并阻止程序正常执行。}
      \item {\textbf{修复并完成闭环验证}：修复 MEM 镜像生成逻辑，确认工具链生成的 MEM 镜像从错误的 5887 Words 恢复至预期的 5888 Words，重新加载后完成 FPGA 验证；成功观测到目标地址的 13 次字符写入和预期输出结果，并进一步编写 C 程序算法完成硬件原型验证，验证修复后的软件-工具链-SoC-FPGA 执行闭环。}
    \end{cvitems}
  }

  %---------------------------------------------------------
  \cventry
  {核心开发者} % Role
  {\textbf{TimeArc} - 本地优先的跨平台时间追踪与分析应用} % Project
  {开源合作项目} % Course/lab/team/personal project
  {2026 年 3 月 - 至今} % Date(s)
  {
    \begin{cvitems} % Description(s) of goals/contributions/results
      \item {\textbf{项目目标}：构建本地优先的时间追踪工具，自动记录前台应用、窗口及后台音视频活动，并提供分时段、分应用或分类统计等功能。}
      \item {\textbf{架构与存储接口}：设计 GUI 和后台数据采集服务两个进程协作的架构；明确 GUI 与后台服务的数据库权限，将采集历史与界面设置数据分离；采用 C 语言设计存储层的统一接口，支持 macOS 通过 Swift 桥接实现跨语言调用。}
      \item {\textbf{原生数据采集}：选用 Swift 语言结合 macOS API，将采集逻辑拆分为探针、协调器与状态机实现；通过 Protocol 解耦系统 API 和业务逻辑，处理应用切换、空闲状态、媒体信息及数据落库，实现前台应用、窗口标题与音视频活动探测。}
      \item {\textbf{分类规则与统计可视化}：使用 C++ 实现统一的应用与网站分类系统，内置 239 条规则、13 个类别，支持用户编辑、禁用、新建或恢复默认规则；设计钟形分类统计图，实现短片段过滤、活动区间去重以及相邻片段归并。}
      \item {\textbf{客户端与工程交付}：使用 Qt 6 与 Objective-C++ 完善 macOS 原生菜单、状态栏、窗口控制和快捷键交互；统一桌面与移动界面的多语言文案；完成 Bash 构建脚本与 macOS DMG 打包流程。}
    \end{cvitems}
  }

  %---------------------------------------------------------
\end{cventries}
